\thispagestyle{plain}
\addcontentsline{toc}{chapter}{Abstract}
\begin{center}
  {\huge\itshape Abstract\par}
\end{center}
\vspace{0.5em}

\noindent\rule{\linewidth}{0.5mm}
\textbf{Name of the Student}: \ThesisAuthor
\hfill \textbf{Roll No.}: \ThesisRollNumber\\
\textbf{Degree}: \ThesisDegreeShort
\hfill \textbf{Department}: Electrical Engineering\\
\textbf{Thesis Title}: \ThesisTitle\\
\textbf{Thesis Supervisor}: \ThesisSupervisor\\
\textbf{Thesis Co-Supervisor}: \ThesisCoSupervisor\\
\textbf{Month and Year of Submission}: \ThesisSubmissionDate
\par\noindent\rule{\linewidth}{0.5mm}

This thesis investigates upright stabilization of a custom rotary (Furuta)
inverted pendulum driven by a stepper motor. Unlike the torque-controlled
actuators commonly assumed in analytical studies, the experimental platform
accepts open-loop STEP/DIR speed commands while measuring the actual arm and
pendulum angles. The resulting mismatch between commanded and realized motion
is central to the modelling, controller implementation, and interpretation of
the experiments.

A first-principles nonlinear model is derived and linearized about the upright
equilibrium. An acceleration-input formulation is then connected to the
firmware path in which commanded acceleration is integrated into step rate.
Three upright-only controllers are implemented on the same hardware: linear
full-state feedback, a hybrid pendulum-surface sliding-mode controller with
base-centering assistance, and a four-state sliding-mode controller. The
real-time implementation includes filtered derivative estimation, wrap-safe
angle processing, sensor-glitch rejection, reference-profile supervision,
actuator limits, stall diagnostics, and controlled disarming.

Experimental evidence is drawn from seven pinned trials: three undisturbed
hold trials, one linear-controller reference-tracking trial, and three
manually disturbed tap trials. In these representative trials, the linear
controller produced the lowest hold RMS pendulum error and was the only
controller to complete the tap trial without falling. The hybrid
sliding-mode controller produced the smallest peak pendulum excursion in the
hold condition. The four-state controller required substantially greater
command effort and was more sensitive to actuator shaping. Reference motion
was achieved without a fall, but none of the commanded steps satisfied the
fixed settling criterion within the available windows. Because each condition
contains one trial, the comparison is a qualitative engineering case study,
not a statistical controller ranking.
